\documentclass[11pt,a4paper]{article}
\usepackage[utf8]{inputenc}
\usepackage[T1]{fontenc}
\usepackage{lmodern}
\usepackage{amsmath,amssymb,amsfonts,amsthm}
\usepackage{geometry}
\geometry{margin=1in}
\usepackage{xcolor}
\usepackage{hyperref}
\usepackage{cleveref}
\usepackage{booktabs}
\usepackage{array}
\usepackage{microtype}
\usepackage{tcolorbox}

\hypersetup{colorlinks=true, linkcolor=blue, citecolor=red, urlcolor=cyan,
  pdftitle={A Falsification-Ready Method for the K3xT2 Reading of the Cooper s7/s10 Families: Exact Computations, Formal Verification, and Controls},
  pdfauthor={Xavier Callens and the SocrateAI Scientific Agora Collaboration}}

\theoremstyle{remark}
\newtheorem{remark}{Remark}[section]
\newcommand{\lean}[1]{\texttt{#1}}
\newcommand{\tierB}{\colorbox{yellow!20}{\textsf{\scriptsize Tier B --- derived, not measured}}}
\newcommand{\tierL}{\colorbox{blue!10}{\textsf{\scriptsize Tier L --- literature}}}
\newcommand{\tierC}{\colorbox{orange!15}{\textsf{\scriptsize Tier C --- blocked physics}}}

\newtcolorbox{scopebox}{colback=gray!5, colframe=gray!45, boxrule=0.4pt, arc=1pt, left=6pt, right=6pt, top=3pt, bottom=3pt,
  fonttitle=\bfseries\small, title={Scope}}
\newtcolorbox{headlinebox}{colback=red!4, colframe=red!45!black, boxrule=0.5pt, arc=1pt, left=6pt, right=6pt, top=4pt, bottom=4pt,
  fonttitle=\bfseries\small, title={Result}}
\newtcolorbox{correctionbox}{colback=orange!6, colframe=orange!55!black, boxrule=0.5pt, arc=1pt, left=6pt, right=6pt, top=4pt, bottom=4pt,
  fonttitle=\bfseries\small, title={Corrections recorded in band}}

\title{\Large \textbf{A Falsification-Ready Method for the $K3\times T^2$ Reading}\\[0.2em]
\textbf{of the Cooper $s_7$/$s_{10}$ Families:}\\[0.2em]
\textbf{Exact Computations, Formal Verification, and Controls}}
\author{\textbf{Xavier Callens} \and \textbf{SocrateAI Scientific Agora Collaboration} \\[0.3em]
\textit{Stream 2 --- Selection \& Geometry} \\
\textit{(with the simulator stream's ledger, as reported by that stream)}}
\date{2026-10-10 \quad (preprint; not reviewed)}

\begin{document}
\maketitle

\begin{scopebox}
This report is about method and exact mathematics. Its object is the Shioda--Inose reading of $K3\times T^2$ for two families of
$K3$ surfaces. Results are stated at the tier of their evidence. The physical reading of $K3\times T^2$ as a compactification, and any
dark-sector observable, are \tierC\ and are outside what is computed here.
\end{scopebox}

\begin{abstract}
We describe, and exercise on two families, a method for turning a geometric hypothesis into claims that can fail. Each claim is
reduced to a statement an exact checker or a proof kernel can decide; each checker ships controls that must fail on tampered input;
every number in a paper is read from a certificate by a script that refuses on drift; and a failing certificate has a stated
consequence. Applied to the Shioda--Inose reading of $K3\times T^2$ it gives the following. The transcendental lattice of
$E\times E'$ for a cyclic $n$-isogeny is computed exactly and is isometric to $U\oplus\langle 2n\rangle$ at every level $n\le24$ and to
no other $U\oplus\langle 2m\rangle$ (24 of 24 on the diagonal, 0 of 552 off it). It equals the certified lattice of each family at
$n=7$ and $n=10$. The level read from the lattice agrees with the level read from the modular side. At $z=1/27$ the section required
by the $M_7$ polarization is constructed exactly, reconstructed from $52$ specializations by an over-determined fit, and read:
$\bar P\cdot\bar O=5$, contact $0$ with the $A_1$ fibre, height $14$, equal to the lattice value. We list what would falsify each
claim and how the pipeline is re-pointed at a new candidate by changing one integer.
\end{abstract}

\tableofcontents

\section{The approach}\label{sec:approach}

The programme's hypothesis has a mathematical core (a family of $K3$ surfaces with a symmetric-square structure, a lattice, a
modular coordinate) and a physical superstructure. The method keeps the two apart and makes the core falsifiable. It has six parts.

\begin{enumerate}
\item \textbf{Every claim carries its tier.} Kernel-checked statements (Tier A), exact derivations (Tier B), literature (Tier L) and
  blocked physics (Tier C) are never written in one another's language. A prose lint applies the same forbidden-verb and conjecture
  rules to the whole body of each paper, and blocks fibre-type tokens that the programme's ruling forbids.
\item \textbf{Numbers are read, not typed.} Each table is generated from a certificate by one script. Its \texttt{--check} mode exits
  non-zero if a fragment drifts, and its controls tamper with the certificates and require the table to change.
\item \textbf{Exact arithmetic, over-determined.} Lattices, intersection numbers and heights are computed in exact integer or rational
  arithmetic. Where a function must be reconstructed from specializations, the number of points exceeds the number of unknowns, the
  solution must be unique, and further points are held out and tested.
\item \textbf{Controls that must fail.} Every checker ships negative controls: a wrong level, a non-cyclic isogeny, a reversed
  orientation, a tampered lattice, a wrong normalization, too few points, a different section. A control that cannot fail is not a
  control.
\item \textbf{Formal verification where it reaches.} The symmetric-square factorization, the action of $\Gamma_0(N)^+$ on
  $U\oplus\langle2N\rangle$ and the rank-jump lattice rows are kernel-checked in Lean~4 by the theory stream, with a second
  independent development for six of the rows. Release gates (build, absence of \texttt{sorry}, axiom audit, statement lock,
  producer $\ne$ verifier) are run by a session that did not write the code. The recorded rule is that the claim must live in the
  statement: a theorem can be true and compile while proving less than its name says.
\item \textbf{Removal and correction rules.} A failed hard gate removes a candidate; an absent certificate does not; consistency gates
  never score; and any correction is dated and machine-visible, not a footnote.
\end{enumerate}

\begin{table}[h]
\centering\footnotesize
\begin{tabular}{p{5.4cm}p{3.4cm}p{5.4cm}}
\toprule
Result & Tier & Source of the number \\
\midrule
$T(E\times E')\cong U\oplus\langle 2n\rangle$, $n=7,10$ & B (exact computation) & \nolinkurl{K3xT2_READING_S.json} \\
The same, for every $n\le24$ against every target & B & \nolinkurl{K3xT2_READING_S_LEVEL_SWEEP.json} \\
$T(E\times E)$ at the three $s_7$ points of $\rho=20$ & B, forced & \nolinkurl{K3xT2_READING_S.json} \\
Level consistency, lattice $n$ = modular $n$ & B (hard gate) & \nolinkurl{T3_LEVEL_CONSISTENCY.json} \\
The section at $z=1/27$: $\bar P\cdot\bar O$, contact, height & B (exact; construction L) & \nolinkurl{TW2_SECTION_DESCENT.json} \\
Counts of the simulator stream's cross-check ledger & mirrored, not verified here & \nolinkurl{refs/simulator_k3t2_ledger_counts_v3.json} \\
Identification of the $T^2$ of a compactification with $E$ & C, not made & --- \\
\bottomrule
\end{tabular}
\caption{What this report computes and where each number comes from. B: derived, not measured. L: literature. C: blocked physics.}\label{tab:map}
\end{table}

\section{Results}\label{sec:results}

\subsection{The lattice of $E\times E'$, and a sweep over levels}\label{sec:lattice}

Let $E\times E'=\mathbb{C}^2/\mathbb{Z}^4$. Then $H^2(E\times E',\mathbb{Z})\cong\wedge^2\mathbb{Z}^4$ with the pairing given by the
coefficient of $e_1\wedge e_2\wedge e_3\wedge e_4$ in $u\wedge w$, and a two-dimensional subtorus spanned by $v_1,v_2$ has class
$v_1\wedge v_2$. For curves without complex multiplication (an assumption that is stated, not computed), the N\'eron--Severi lattice is
spanned by the two fibres and the graph of the isogeny. The checker verifies signature $(1,2)$ and saturation, computes the
orthogonal complement as an integral basis, and finds an explicit unimodular basis change onto the target.

\begin{table}[h]
\centering\scriptsize\setlength{\tabcolsep}{3pt}
\begin{tabular}{llllll}
\toprule
Family & $n$ & NS sig. & saturated & $T(E\times E')$ (Gram), sig. & $\cong$ certified $T$ \\
\midrule
\lean{cooper\_s7} & $7$ & $(1,2)$ & yes & $[0,0,-1;0,14,0;-1,0,0]$, $(2,1)$ & yes \\
\lean{cooper\_s10} & $10$ & $(1,2)$ & yes & $[0,0,-1;0,20,0;-1,0,0]$, $(2,1)$ & yes \\
\midrule
$s_7$ point & model $j$ & CM disc. & $T(E\times E)$ & certified $T$ & agree \\
\midrule
$1/27$ & $16581375$ & $-28$ & $[2,0;0,14]$ & $[2,0;0,14]$ & yes (forced) \\
$-1$ & $-3375$ & $-7$ & $[2,1;1,4]$ & $[2,1;1,4]$ & yes (forced) \\
$\infty$ & $0$ & $-3$ & $[2,1;1,2]$ & $[2,1;1,2]$ & yes (forced) \\
\bottomrule
\end{tabular}

\caption{Top: the computed $T(E\times E')$ against each family's certified $T$. Bottom: $T(E\times E)$ at the $s_7$ points of Picard
number $20$, with the CM order identified by computation from the model's $j$. The bottom agreement is forced (an elliptic point of
the modular curve is a CM point), so it is a consistency check.}\label{tab:readings}
\end{table}

\begin{table}[h]
\centering\small
\begin{tabular}{lr}
\toprule
Level sweep (exact; assumes no CM) & Result \\
\midrule
levels run & $1,\dots,24$ \\
(computed lattice, offered target) pairs & $576$ \\
isometric on the diagonal (same level) & $24$ of $24$ \\
isometric off the diagonal (wrong level) & $0$ of $552$ \\
N\'eron--Severi of signature $(1,2)$ and saturated at every level & yes \\
\bottomrule
\end{tabular}

\caption{The same code at every level: each computed lattice is compared with every target $U\oplus\langle2m\rangle$. The isometry
exists exactly when $m=n$.}\label{tab:sweep}
\end{table}

The sweep is the evidence that the two carried levels, $7$ and $10$, are not special cases: a wrong level is rejected by the code that
accepts the right one, at every one of the 576 pairs.

\subsection{Level consistency}\label{sec:t3}

\begin{table}[h]
\centering\small
\begin{tabular}{llll}
\toprule
Family & $n$ from lattice & $n$ from modular side & Verdict \\
\midrule
\lean{cooper\_s7} & $7$ & $7$ & T3\_AGREE(n=7) \\
\lean{cooper\_s10} & $10$ & $10$ & T3\_AGREE(n=10) \\
\bottomrule
\end{tabular}

\caption{The level $n$ read from the lattice certificate against the level read from the modular side (an eta-quotient Hauptmodul
with Newman's conditions and an Atkin--Lehner check). Both legs start from the same operator, so agreement is a consistency check and
not an independent measurement.}\label{tab:t3}
\end{table}

\subsection{The section at $z=1/27$}\label{sec:section}

\begin{table}[h]
\centering\small
\begin{tabular}{lrrlrll}
\toprule
$s_7$ member & $|\mathrm{disc}\,T|$ & $\rho$ & extra roots & MW rank & $h(P)$ & $\bar P\cdot\bar O$ \\
\midrule
generic & $14$ & $19$ & --- & $1$ & $14$ & $5$ \\
$z=1/27$ & $28$ & $20$ & $A_{1}$ & $1$ & $14$ & $5$ \\
$z=-1$ & $7$ & $20$ & $A_{1}$ & $1$ & $7/2$ & $0$ \\
$z=\infty$ & $3$ & $20$ & $A_{2}$ & $0$ & --- & no section \\
\bottomrule
\end{tabular}

\caption{The generic member and the three $\rho=20$ points of $s_7$ (lattice arithmetic and exact reduction of the model). At
$z=\infty$ the Mordell--Weil rank is $0$.}\label{tab:tw2}
\end{table}

For $z=1/27$ the section is the one attached to the degree-$7$ endomorphism $\sqrt{-7}$ of $E$. It is constructed, not inferred:
\begin{enumerate}
\item The V\'elu isogeny $\varphi\colon E_1\to E_2$ of degree $7$ is written over $\mathbb{Q}$, and $\varphi_y=\varphi_x'$ is checked
  exactly against the curve equation.
\item The cubic $C_t$ of Kumar--Kuwata (3.1) is reduced through $s=x_2-t^2x_1$ to a quartic with a rational point and mapped to a
  Weierstrass model by Mordell's formulas, verified by identity on the actual divisor points and by the quartic's own invariants.
\item The nine-point divisors $D^{\pm}$ ($\varphi_y(x_1)=\pm t^3$) are summed by Riemann--Roch: the function in $L(10\,O)$ vanishing on
  $D^{\pm}$ has one further zero $R$, and $Q^{\pm}=-R$. The section is $P=Q^{+}-Q^{-}$.
\item The model is matched to the simple normal form $Y^2=X^3-3\alpha X+t_s+t_s^{-1}-2\beta$, with $t_s=t^6/343$ and
  $343=\sqrt{\Delta_{E_1}/\Delta_{E_2}}$ computed; the $j$-match and the twist constant are asserted at every specialization.
\item $X'(t_s)$ is reconstructed as a rational function and read.
\end{enumerate}

\begin{table}[h]
\centering\small
\begin{tabular}{p{5.6cm}p{8.6cm}}
\toprule
Quantity & Value \\
\midrule
fit points / unknowns / held out & $44$ / $26$ / $8$ \\
nullspace dimension (1 = unique) & $1$ \\
$\deg N$, $\deg D^2$ of $X'=N/D^2$ & $14$, $10$ \\
held-out points agree & yes \\
$\bar P\cdot\bar O$ (read) & $5$ \\
contact with the $A_1$ fibre, contr (read) & $0$ \\
height $h=2\chi+2\bar P\cdot\bar O-\mathrm{contr}$ (read) & $14$ \\
earlier steps and literature, read from their certificates & step 0: $h=14$; step 2a: $\bar P\cdot\bar O=5$, contr $=0$; Kumar--Kuwata: $h=14$ \\
discriminant orders (model; fibration certificate) & $\{10,10,2,1,1\}$ ; $\{10,10,2,1,1\}$ \\
cross-check agrees & yes \\
\bottomrule
\end{tabular}

\caption{The reconstruction and what is read from it. The cross-check row compares with the lattice steps and the fibration
certificate, each read from its own file.}\label{tab:descent}
\end{table}

\begin{headlinebox}
The section meets $O$ five times: once inside the $A_1$ fibre (so on its identity component) and four times at the roots of an
irreducible quartic. The two $E_8$ fibres carry no contact. Its height is $14=2\cdot7$, equal to the value from the lattice arithmetic
and from the Kumar--Kuwata theorem. \tierB.
\end{headlinebox}

\begin{remark}
A control found a structural fact: $\bar P\cdot\bar O=(\deg N-4)/2$ for any denominator, since poles that leave the finite line
reappear at infinity. The count is fixed by the degree of the numerator; the pole locations carry the contact information. Doubling
the section does not fit the same shape, as expected from its height.
\end{remark}

The section is a height-$14$ element; whether it generates the full Mordell--Weil group is the one question left open by this
computation (no index is computed; the lattice and the literature both give $14$).

\subsection{The simulator stream's ledger}\label{sec:sim}

The DualScaleSimulator stream compares, for $332$ sealed targets of the Lean library, values from a separately written pipeline with
the values stated in Lean. We mirror its counts, recounted from its row data by a script that refuses on mismatch and pinned to the
source file's hash. We report them as that stream's record; the rows are not re-derived here.

\begin{table}[h]
\centering\small
\begin{tabular}{lr}
\toprule
Count in the simulator stream's ledger & Value \\
\midrule
sealed targets compared & 332 \\
\quad AGREE & 68 \\
\quad SAME\_DECLARED\_INPUT (downgraded) & 3 \\
\quad NOT\_COMPUTED & 261 \\
\quad DISAGREE & 0 \\
rows with two or more independent routes & 6 \\
declared inputs (of which typed from memory) & 80 (20) \\
rigidity entries: confirmed / disputed / unreviewed & 3 / 2 / 28 \\
chain links consistent; independent corroborations & 8 (all); 3 \\
items its authors could not do & 30 \\
\bottomrule
\end{tabular}

\caption{The simulator stream's ledger, mirrored (source commit and hash pinned in \texttt{refs/}). ``Rigid'' there means a selecting
condition that never mentions the target value was scanned for one parameter.}\label{tab:ledger}
\end{table}

Most targets are not computed, so the ledger is a map of what remains. A fifth of the inputs it declares were typed from memory, which
the ledger records. The same stream's note records that its implemented model sits at level $12$ through a hard-coded constant, that
under the potential actually integrated no observable depends on the level, that an opt-in modular potential does read the level, and
that its level experiment (item~L5) ran only partially.

\section{Falsification and re-adaptation}\label{sec:falsify}

The value of the method is that a change of hypothesis is a change of input, and that each claim has a named way to fail.

\begin{table}[h]
\centering\footnotesize\renewcommand{\arraystretch}{1.35}
\begin{tabular}{p{4.2cm}p{4.4cm}p{5.4cm}}
\toprule
Claim & Would be falsified by & Exercised by \\
\midrule
$T(E\times E')\cong U\oplus\langle2n\rangle$ at level $n$ & a level where the computed lattice is not isometric to its target, or is isometric to a different target & the 24-level sweep: 576 pairs, diagonal only; controls: wrong $n$, non-cyclic isogeny, reversed orientation, tampered Gram \\
A family's certified lattice has the level its modular coordinate has & a certificate with lattice $n\ne$ modular $n$ (removes the candidate) & four other register entries fail the modular leg; 25 controls \\
The $M_7$ polarization needs a section of height $2n$ with $\bar P\cdot\bar O=n-2$ & no such section, or one of a different height, at a locus where the lattice allows it & the explicit construction at $z=1/27$; the rank-$0$ locus $z=\infty$ has none, as the lattice step says \\
An isogeny of degree $d$ gives a section of height $2d$ & a graph whose reconstructed height differs & the doubled section fails the shape; 14 controls \\
A lattice or rank-jump statement of the theory & a failing gate of the Lean release run, or a statement that omits what its name claims & kernel build, axiom audit, statement lock, producer $\ne$ verifier \\
\bottomrule
\end{tabular}
\caption{What would falsify each claim, and what already exercises it.}\label{tab:falsify}
\end{table}

\paragraph{Re-pointing the pipeline.} A new candidate needs: its certified lattice (from the lattice checker); its level from the
modular side; the level-consistency gate; and, if its lattice carries a Picard-number-$20$ point, the section construction at that
point, which takes the isogeny degree and the curve as input. The sweep table reports that the lattice step runs unchanged at every level tried.
A failed gate has a stated consequence (removal, or a dated correction), and every table in a new paper is regenerated by the same
renderer, whose controls fail if a number is typed.

\paragraph{What the method can and cannot decide.} It decides the mathematical core. It does not by itself connect that core to
physics: that connection needs a worked effective-theory matching, a quantitative prediction, a pre-registered pin and a test against
public data, and none of those links exists yet. The method is the same one that would be applied to them.

\section{Corrections recorded in band}\label{sec:corrections}

\begin{correctionbox}
\begin{enumerate}
\item The reading of ``$K3\times T^2$'' used here was first described as the programme's choice before it had been ruled. It was
  ruled on 2026-10-08 (decision D29$'$, taken on the decision authority's behalf under explicit delegation) and is a convention for
  selection criteria only.
\item The first reconstruction of the section used a transcription of the Weierstrass form from a garbled line of the source. The
  exact $j$-match test refused it, and the source's simple normal form was used; the constant $t_s=t^6/343$ was calibrated on three
  points and is asserted at all fifty-two.
\item A first sensitivity control for the read-off passed for the wrong reason (removing the quartic moves its poles to infinity,
  leaving the total unchanged). It was replaced by one that adds a genuine pole.
\item A note elsewhere in the programme said the host lacked a numerical library and that two test modules were ignored. That was a
  property of the system interpreter; under the project's own environment the full suite runs.
\item One continuous-integration run of the simulator failed once on a $10^{-12}$ tolerance in an ODE-solve comparison and passed on
  re-run of the same commit; the tolerance is hardware-sensitive and was not changed.
\end{enumerate}
\end{correctionbox}

\section{Reproducibility}

Every number in \cref{tab:readings,tab:sweep,tab:t3,tab:tw2,tab:descent,tab:ledger} is read from a certificate or a pinned mirror by
\nolinkurl{scripts/render_paper_tables.py}. Checkers and their controls:
\begin{itemize}\small
\item \nolinkurl{checkers/check_K3xT2_reading_S.py} (9 controls);
\item \nolinkurl{checkers/check_K3xT2_reading_S_level_sweep.py} (6 checks);
\item \nolinkurl{checkers/check_T3_level_consistency.py} (25 controls);
\item \nolinkurl{checkers/check_TW2_section_descent.py} (14 controls);
\item \nolinkurl{scripts/extract_simulator_ledger_counts.py} (recount; 8 controls).
\end{itemize}
Decision records: \nolinkurl{briefs/T0_DECISIONS_2026_10_07_STREAM2_DELEGATED.md} (D29$'$--D31$'$);
Lean release-gate record: Stream~1 \nolinkurl{briefs/STREAM2_REVERIFICATION_2026_10_08.md}.

\section*{Provenance}
Generated-by: Claude (Fable~5.1 and Sonnet~5.5), Stream~2 \textbar\ Verified-by: \texttt{render\_paper\_tables.py --check} and its
controls; each cited checker's own controls \textbar\ Reviewed-by: N (preprint; not reviewed).

\begin{thebibliography}{9}
\bibitem{KK} A.~Kumar and M.~Kuwata, ``Elliptic $K3$ surfaces associated with the product of two elliptic curves,''
  \textit{arXiv:1409.2931}.
\bibitem{SS} M.~Sch\"utt and T.~Shioda, ``Elliptic surfaces,'' \textit{arXiv:0907.0298}.
\bibitem{Takatsu} T.~Takatsu, ``On the geometry of singular $K3$ surfaces with discriminant 3, 4 and 7,'' \textit{arXiv:1903.03054}.
\bibitem{Dolgachev} I.~Dolgachev, ``Mirror symmetry for lattice polarized $K3$ surfaces,'' \textit{arXiv:alg-geom/9502005}.
\bibitem{Prev} X.~Callens and the SocrateAI Scientific Agora Collaboration, ``The selected $K3$ surfaces of the Cooper $s_7$/$s_{10}$
  families: identification, the $K3\times T^2$ reading, and the twisted-Weierstrass route,'' Stream~2 note, 2026-10-08.
\end{thebibliography}

\end{document}
